\section{Laurent coefficients of every original test-configuration section}

\begin{theorem}\label{thm:native-tc-global-laurent-sections}
\lean{EnvelopeProperty.NativeTestConfiguration.TestConfiguration.productTotalMap, EnvelopeProperty.NativeTestConfiguration.TestConfiguration.productTotalMap_gaussPoint, EnvelopeProperty.NativeTestConfiguration.TestConfiguration.productTotalLinePowerIso, EnvelopeProperty.NativeTestConfiguration.TestConfiguration.globalProductPowerSection, EnvelopeProperty.NativeTestConfiguration.TestConfiguration.globalPowerLaurentCoefficients, EnvelopeProperty.NativeTestConfiguration.TestConfiguration.globalProductPowerSection_eq_sum}
\leanok
Let $\mathcal X$ be an original test configuration, $M$ its original
linearized module, and $L$ its fibre-one module. For every $n\geq0$,
the native composite of the noncentral product inverse and open inclusion,
together with its source-derived whole line-power isomorphism, sends every original global section of $M^{\otimes n}$
to an actual whole section of the product pullback of $L^{\otimes n}$.
If $L$ is invertible, this section has unique finitely supported
Laurent coefficients in the entire original space
$H^0(X,L^{\otimes n})$ and is their actual monomial sum. All compactness,
separatedness and quasicoherence needed here follow from the original
test-configuration data and invertibility, rather than being additional
comparison inputs.
\uses{thm:native-test-configuration-whole-open-sections,thm:fc-scheme-module-laurent-coefficients}
\end{theorem}
\begin{proof}
The original fibre-one isomorphism and proper test-configuration
projection imply properness of the original $X$ over the field.
Properness gives compactness and separatedness. Every tensor power
of an invertible module is invertible and hence quasicoherent.
Compose the actual product inverse with the original noncentral open
inclusion to obtain a morphism $g$ from the product to the total space.
The native pullback-composition isomorphisms and the source-derived
whole noncentral line-power isomorphism give an actual isomorphism
$g^*M^{\otimes n}\simeq\operatorname{pr}_X^*L^{\otimes n}$.
Send each original section by the genuine pullback-adjunction unit
along $g$, then by this actual whole module isomorphism. Composition
of the native valued-point maps shows that $g$ sends the whole
product Gauss point to the actual total-space Gauss point. Apply the whole Laurent
coefficient equivalence to this actual product section. Its inverse
formula is the finite sum of original monomial multiples of actual
pullback-unit original sections, which proves both reconstruction
and uniqueness. No homogeneous-coordinate lifting or section
matching property is supplied.
\end{proof}

\begin{theorem}\label{thm:native-tc-global-laurent-fixed-family}
\lean{EnvelopeProperty.NativeTestConfiguration.TestConfiguration.globalProductPowerFamily, EnvelopeProperty.NativeTestConfiguration.TestConfiguration.globalPowerLaurentTerms, EnvelopeProperty.NativeTestConfiguration.TestConfiguration.globalPowerLaurentCoefficient_mem_terms}
\leanok
For every finite family $G$ of original total-space sections in a
fixed degree, the resulting original coefficient/weight family is
finite and is chosen before all original valued points, radii,
local frames and evaluations. Every coefficient of every transported
section is supported in that same original family. The construction
includes degree zero, zero coefficients and negative integer weights.
\uses{thm:native-tc-global-laurent-sections}
\end{theorem}
\begin{proof}
Map the actual finite family by the original restriction and whole
product-section map. Take the finite union of the supports of the
finitely supported coefficient families, retaining each coefficient
with its integer exponent. A coefficient in one of these supports
belongs to this finite union by its defining membership witness.
The construction uses no valued point, radius, frame or functional.
All integer exponents and all original sections are retained.
\end{proof}
